\section{The exact local recurrence and its holomorphic asymptotic}
\label{app:selected-source}

This section proves the local source estimate used in the completed-root
theorem. Its hypothesis is a fixed compact interval
\begin{equation}\label{eq:b1-high-interval}
 I=[\tau_0,\tau_1]\subset(1/4,1/2)
\end{equation}
for every retained complete leaf coordinate. All constants may depend on
this interval and the fixed source tree. The color neighborhood and every
complex index neighborhood below are fixed independently of the level.
We work with the actual finite low core and its specified factorial
continuation. No homogeneous solution of a recurrence is substituted for
that source. In particular, an inverse-power approximation on the real
index grid will not be used as if it were an exponentially accurate
complex interpolation.

\subsection{Canonical factorials and the literal mixed coefficient}

To move a color parameter off the integer grid, we must specify a function that agrees with the finite coefficient and preserves its phase. The Gamma product below supplies that function for the sine factorials. Its exact shift identity will then allow the finite recurrence to remain valid at complex colors; agreement only in a leading asymptotic term would not be sufficient.

Retain
\[
 d=N-1,\quad M=\lfloor d/2\rfloor,\quad
 q=e^{\pi\ii/N},\quad Q=q^2,\quad \kappa=q^{d^2/2},\quad
 F_r=(Q;Q)_r.
\]
In this section $x$ denotes a scaled Newton index, rather than the
physical-frame scalar $q^d$. The sine factorial has the exact continuation
\begin{equation}\label{eq:b1-factorial}
 G_N(z)=\left(\frac{2\pi}{N}\right)^z\Gamma(z+1)
 \prod_{k=1}^{\infty}
 \frac{\Gamma(Nk+z+1)}{\Gamma(Nk-z)(Nk)^{2z+1}},
 \qquad -1<\Re z<N.
\end{equation}
All positive-base powers use their real logarithms. The logarithm of the
$k$th factor is $O_K(k^{-2})$ on compact subsets of the strip: in the
Gamma-ratio expansion its $k^{-1}$ term cancels. Thus the product is
holomorphic and zero-free there. Gamma's functional equation and the
convergent sine product give
\begin{equation}\label{eq:b1-factorial-shift}
 \frac{G_N(z+1)}{G_N(z)}=2\sin\frac{\pi(z+1)}N,
 \qquad G_N(0)=1,\qquad G_N(r)=s_r\quad(0\le r<N).
\end{equation}
The marked quantum factorial is
\begin{equation}\label{eq:b1-marked-factorial}
 \widehat F_N(z)=
 e^{\pi\ii z(z+1)/(2N)-\pi\ii z/2}G_N(z).
\end{equation}
Consequently
\[
 \widehat F_N(z+1)=(1-Q^{z+1})\widehat F_N(z),
 \qquad \widehat F_N(r)=F_r,
 \qquad Q^z=e^{2\pi\ii z/N}.
\]
These identities fix the continuation and its phase, not just its
asymptotic expansion. On a compact subset of $0<\Re z<1$, direct
Stirling expansion of \eqref{eq:b1-factorial} gives
\begin{equation}\label{eq:b1-factorial-asymptotic}
 \begin{aligned}
 \log G_N(Nz)
 &=NJ(z)+\frac12\operatorname{Log}(2N\sin\pi z)
   +\frac{\pi\cot\pi z}{12N}+O(N^{-3}),\\
 J(z)&=\int_0^z\operatorname{Log}(2\sin\pi u)\,\dd u.
 \end{aligned}
\end{equation}
Here the logarithm continues the real positive factorial branch. One
way to verify the uniform remainder is to expand the $k$th Gamma ratio:
its error is $O_K(N^{-3}k^{-3})$. These errors are summable; the
leading and half-log terms sum by the sine product. The first correction
sums by
\[
 z^{-1}+\sum_{k\ge1}\bigl((k+z)^{-1}-(k-z)^{-1}\bigr)
 =\pi\cot\pi z.
\]
Cauchy's estimate supplies the corresponding derivative bounds on
smaller compact sets.

Let $h$ be a physical color deficit and put $t=h/N$. At an integer
index $r>h$ the pure core coefficient is the literal factorial expression
\begin{equation}\label{eq:b1-pure-coefficient}
 A_N(r,t)=\frac{N^2s_{2r}}
 {s_r^2G_N(r-Nt)G_N(r+Nt)}.
\end{equation}
For a retained half-twist sign $\eps\in\{1,-1\}$, the exact triangular
Newton transport gives
\begin{equation}\label{eq:b1-signed-finite}
 d_N^\eps(r,t)=\kappa^\eps
 \sum_{v=r}^{M}(-1)^v q^{\eps v(r+1)}
 \frac{s_v}{s_rs_{v-r}} A_N(v,t),
 \qquad d_N^0(r,t)=A_N(r,t).
\end{equation}
For complex $t$ in a sufficiently small fixed disk, only the two
color-shifted factorials are continued. All their arguments remain in
regular strips. Equation \eqref{eq:b1-signed-finite} is then an exact
finite, even holomorphic function. At every physical integer $h<r$ it
is the original transported core, including all first normal directions.

The cancellation inside this finite sum is encoded by
\begin{align}
 A_{h,r}
 &=\sum_{v=r}^{M}
 \frac{Q^{-vr}F_{2v}}
 {F_v\widehat F_N(v-h)\widehat F_N(v+h)F_{v-r}},
 \label{eq:b1-A-source}\\
 D_{h,r}
 &=\frac{Q^{-r^2}F_{2r}}
 {F_r\widehat F_N(r-h)\widehat F_N(r+h)},
 \qquad U_{h,r}=A_{h,r}/D_{h,r}.
 \label{eq:b1-D-source}
\end{align}
The sum starts at $r$ in the present corridor. For physical $r<h$ its
lower endpoint would be $h$; those low indices belong to the separate
global envelope, not to the interpolation asserted here. The terminal
source in \eqref{eq:b1-D-source} is exactly $U_{h,M}=1$.

\subsection{Deriving the recurrence, including its endpoint}

The identity traditionally called a Ward equation here is simply a second-order recurrence obtained by telescoping the actual finite sum. A recurrence has many solutions, so we include its upper endpoint condition in the derivation. That condition is what selects the coefficient occurring in the invariant.

Put $H=Q^h=e^{2\pi\ii t}$, $T_h=H+H^{-1}$ and $Y=Q^r$. The
literal finite sum satisfies
\begin{multline}\label{eq:b1-Ward}
 QY^2(1-Q^2Y)A_{h,r+2}
 +Y(T_hQ^3Y^2+Q+1)A_{h,r+1}\\
 -Q(QY+1)(QY^2-1)A_{h,r}=0,
 \qquad A_{h,M+1}=0.
\end{multline}
Here and below an absent term beyond the finite upper support is zero.
We give the summation certificate, since the endpoint fixes which
solution of this second-order equation is the physical one.

For the summand of \eqref{eq:b1-A-source}, write $z=Q^v$. Its ratio
under $v\mapsto v+1$ is
\[
 \frac{Y^{-1}(1-Qz^2)(1+Qz)}
 {(1-Qz/H)(1-QHz)(1-Qz/Y)}.
\]
Use the rational certificate
\[
 g_h(z)=\frac{(1-z/H)(1-Hz)(1-z/Y)}{z^2}.
\]
The telescoping difference, divided by the original summand, is
\begin{equation}\label{eq:b1-Ward-certificate}
 \left(\frac1{YQ^2}-1\right)z^{-2}
 +\left(\frac1{YQ}+H+H^{-1}+\frac1Y\right)z^{-1}
 -\frac1{YQ}-1-\frac{H+H^{-1}}Y.
\end{equation}
The lower endpoint has $z=Y$ and hence vanishes. The upper endpoint
contains $F_{2M+2}=0$. The ratio for the shift $r\mapsto r+1$ is
$z^{-1}-Y^{-1}$. Expressing the first two $z^{-1}$ moments in terms of
$A_{h,r+1}$ and $A_{h,r+2}$ in
\eqref{eq:b1-Ward-certificate} gives \eqref{eq:b1-Ward} by multiplication
of its displayed denominators. This argument remains exact for complex
$h=Nt$: the shifted factorial equations are exact and the lower zero
is still $z=Y$. It does not use interpolation from finitely many colors.

Normalize by the actual $D_{h,r}$ and set
\[
 \operatorname{den}_r=(1-QY/H)(1-QYH).
\]
Then
\begin{equation}\label{eq:b1-normalized-recurrence}
 U_r=p_N(r,t)U_{r+1}+q_N(r,t)U_{r+2},
\end{equation}
where the exact rational coefficients are
\begin{align}
 p_N&=-\frac{Q^{-2}Y^{-1}(T_hQ^3Y^2+Q+1)}
                 {\operatorname{den}_r},
 \label{eq:b1-p-coefficient}\\
 q_N&=-\frac{Q^{-4}Y^{-2}(1-Q^2Y)(1-Q^3Y^2)(1+Q^2Y)}
 {\operatorname{den}_r\operatorname{den}_{r+1}}.
 \label{eq:b1-q-coefficient}
\end{align}
All these denominators are bounded away from zero on the corridor used
below. At the last recurrence row one has the exact wall
\begin{equation}\label{eq:b1-wall}
 q_N(M-1,t)=0,
 \qquad U_{M-1}=p_N(M-1,t),\qquad U_M=1.
\end{equation}
For odd $N=2m+1$, $M=m$ and the factor $1-Q^3Y^2$ vanishes.
For even $N=2m$, $M=m-1$ and $1+Q^2Y$ vanishes.
The zero is independent of $H$. In particular no nonexistent
$D_{h,M+1}$ is divided out and no arbitrary terminal ratio is imposed.

\subsection{The selected analytic root and a source-exact continued fraction}

We solve the recurrence backwards from its finite upper endpoint. On the strict high-index interval, one characteristic root dominates the other. The resulting contraction permits a holomorphic continued fraction. Crucially, an exact zero in the last recurrence row erases the artificial tail of this continued fraction, so it agrees with the original finite solution on the integer grid.

Regard \eqref{eq:b1-p-coefficient}--\eqref{eq:b1-q-coefficient} as
holomorphic functions $p(x,t,\varepsilon_N)$ and
$q_c(x,t,\varepsilon_N)$, where
$\varepsilon_N=N^{-1}$, $Y=e^{2\pi\ii x}$ and
$Q=e^{2\pi\ii\varepsilon_N}$. At $\varepsilon_N=0$, put
\begin{equation}\label{eq:b1-leading-recurrence}
 \begin{split}
 D_t&=(1-Y/H)(1-YH),\\
 p_0&=-Y^{-1}\frac{(H+H^{-1})Y^2+2}{D_t},\qquad
 q_0=-Y^{-2}\frac{(1-Y^2)^2}{D_t^2}.
 \end{split}
\end{equation}
The two roots of $\mu^2-p_0\mu-q_0=0$ at $t=0$ are
\begin{equation}\label{eq:b1-two-roots}
 \mu_s(x,0)=-Y^{-1},\qquad
 \mu_f(x,0)=Y^{-1}\cot^2(\pi x).
\end{equation}
On a compact real interval in $(1/4,3/4)$ containing
$[\tau_0,1/2]$, the selected branch $\mu=\mu_s$ is nonzero,
$\Re\mu>0$, and $|\mu_f/\mu|<1$. The same inequalities, with uniform
strict constants, hold on a fixed complex neighborhood and a sufficiently
small complex $t$ disk. Indeed at $t=0$ these assertions follow directly
from \eqref{eq:b1-two-roots}; nonvanishing and a separated quadratic root
persist on a small complex thickening. Write
\begin{equation}\label{eq:b1-dominance}
 \left|\frac{\mu_f}{\mu}\right|\le k_0<1,
 \qquad \Re\mu\ge m_0>0.
\end{equation}
The endpoint $x=1/2$ is included. Direct substitution gives
$\mu(1/2,t)=1$, while $q_0$ has a double zero there and its first
$\varepsilon_N$ derivative also vanishes. The selected root and its
normalization are regular at the physical terminal point.

Choose $\sigma_0>0$ so small that
$[\tau_0-\sigma_0,1-\tau_0+4\sigma_0]\subset(1/4,3/4)$.
Use a fixed complex neighborhood of this larger interval on which
\eqref{eq:b1-dominance} holds, and put
\[
 \mathcal L_N=\left\lceil(1/2-\tau_0+2\sigma_0)N\right\rceil,
 \qquad \mathcal F_x(z)=p(x,t,\varepsilon_N)
                  +\frac{q_c(x,t,\varepsilon_N)}z.
\]
Define the finite continued fraction
\begin{equation}\label{eq:b1-ratio-interpolant}
 v_N(x,t)=\mathcal F_x\circ\mathcal F_{x+\varepsilon_N}\circ
 \cdots\circ\mathcal F_{x+(\mathcal L_N-1)\varepsilon_N}
 \bigl(\mu(x+\mathcal L_N\varepsilon_N,t)\bigr).
\end{equation}
This is defined on a fixed complex neighborhood of $[\tau_0,1/2]$.
For a constant $C$ independent of $N$, the disks of radius
$C\varepsilon_N$ around the respective selected roots are invariant
under backward iteration. To prove this, note first that
\[
 \mathcal F_x(\mu(x+\varepsilon_N,t))-\mu(x,t)
 =O(\varepsilon_N).
\]
On those disks the derivative in the continued-fraction argument has
modulus at most $k_0+O(\varepsilon_N)$. Choose $k_1\in(k_0,1)$ and then
$C$ so that the $O(\varepsilon_N)$ error plus
$k_1C\varepsilon_N$ is at most $C\varepsilon_N$. This proves the
invariance inductively. The disks avoid zero, so every iterate is
holomorphic, nonzero, and has uniformly bounded inverse. In particular
\begin{equation}\label{eq:b1-ratio-leading}
 v_N(x,t)=\mu(x,t)+O(N^{-1})
\end{equation}
uniformly holomorphically on a fixed smaller domain.

At a physical $x=r/N$, the continued fraction reaches the row $M-1$.
By \eqref{eq:b1-wall}, that row forgets its entire remaining tail and
takes exactly the value $p_N(M-1,t)$. Backward induction from the literal
source $U_M=1$ therefore proves
\begin{equation}\label{eq:b1-exact-ratio}
 v_N(r/N,t)=\frac{U_{Nt,r}}{U_{Nt,r+1}},
 \qquad \tau_0N\le r<M.
\end{equation}
This also proves that the displayed physical source values are nonzero.
The equality holds uniformly for complex $t$ in the fixed disk. There is
no assertion about a nonexistent physical ratio at $r=M$.

Adding one more future step to \eqref{eq:b1-ratio-interpolant} changes
the future value by $O(\varepsilon_N)$, and the preceding
$\mathcal L_N$ contractions reduce that change by $k_1^{\mathcal L_N}$.
Thus its Riccati defect is exponentially small. This explains why the
analytic tail used to define the interpolant cannot hide a macroscopic
change of the original source.

\subsection{First transport correction and absolute amplitude normalization}

Let $p_1=\partial_{\varepsilon}p|_{\varepsilon=0}$ and
$q_1=\partial_{\varepsilon}q_c|_{\varepsilon=0}$, differentiating the
explicit rational coefficients. Define
\begin{equation}\label{eq:b1-nu}
 \nu=\frac{p_1+q_1/\mu-q_0\mu_x/\mu^2}
                  {1+q_0/\mu^2}.
\end{equation}
The denominator is $1-\mu_f/\mu$, bounded away from zero by
\eqref{eq:b1-dominance}. Substitution of
$\mu+\varepsilon_N\nu$ into the one-step Riccati equation gives a
residual $O(\varepsilon_N^2)$. The same contracting backward iteration
then gives
\begin{equation}\label{eq:b1-ratio-first}
 v_N=\mu+N^{-1}\nu+O(N^{-2})
\end{equation}
on a fixed smaller complex neighborhood. For example, subtract the
trial iterate from the exact one at each step. The resulting error is
bounded by the geometric sum of its $O(N^{-2})$ local defects and the
contracted terminal discrepancy. Cauchy estimates give the corresponding
uniform derivative expansions.

Use the logarithm with positive real part in
\eqref{eq:b1-dominance}, continuing from $\mu(1/2,t)=1$, and set
\begin{equation}\label{eq:b1-action-amplitude}
 \begin{split}
 S(x,t)&=\int_x^{1/2}\operatorname{Log}\mu(u,t)\,\dd u,\\
 C(x,t)&=\sqrt{\mu(x,t)}
 \exp\left(\int_x^{1/2}\frac{\nu(u,t)}{\mu(u,t)}\,\dd u\right).
 \end{split}
\end{equation}
The square root is one at $1/2$. The integral branches are fixed by the
actual selected source; $C$ is holomorphic and nonzero.

Summing logarithms of the exact ratios in
\eqref{eq:b1-exact-ratio}, using \eqref{eq:b1-ratio-first} and the first
Euler--Maclaurin endpoint term, gives
\begin{equation}\label{eq:b1-U-leading}
 U_{Nt,r}=C(r/N,t)e^{NS(r/N,t)}[1+O(N^{-1})].
\end{equation}
The product stops at $M-1$, because $U_M=1$. Replacing the true endpoint
$M/N$ by $1/2$ costs only $O(N^{-1})$: the displacement is $O(N^{-1})$
at either parity, and $\operatorname{Log}\mu(1/2,t)=0$. The remaining
endpoint terms in the correction are likewise $O(N^{-1})$. In
particular there is no free terminal matching constant.

At $t=0$, direct differentiation of
\eqref{eq:b1-p-coefficient}--\eqref{eq:b1-q-coefficient} yields
\[
 \operatorname{Log}\mu=\pi\ii-2\pi\ii x,
 \qquad \frac\nu\mu=-\frac\pi2\cot(\pi x)-\frac{5\pi\ii}2.
\]
Therefore
\begin{equation}\label{eq:b1-complete-calibration}
 S(x,0)=\pi\ii(1/2-x)^2,\qquad
 C(x,0)=\sqrt{\sin\pi x}\,
                 e^{-3\pi\ii(1/2-x)/2}.
\end{equation}
This calculation fixes the absolute amplitude from the finite source,
rather than from the value of an unknown closed invariant.

\subsection{Exponential interpolation before an oscillatory saddle}

Equation \eqref{eq:b1-U-leading} is not sufficient by itself for a
complex wire saddle. Its relative error is only inverse polynomial and
could exceed a fixed oscillatory cancellation. The exact ratio
\eqref{eq:b1-exact-ratio} supplies the needed stronger estimate.

Put $f_N(x,t)=\operatorname{Log}v_N(x,t)$ and $b_N=M/N$.
On a fixed larger complex tube the $f_N$ are uniformly bounded and
holomorphic. On an inner tube, Cauchy's estimate gives constants
$B,a>0$, independent of $N,t$, such that
\begin{equation}\label{eq:b1-Cauchy}
 |\partial_x^k f_N(x,t)|\le Bk!a^{-k}.
\end{equation}
For $K_N=\lfloor c_aN\rfloor$, with a sufficiently small fixed
$c_a>0$, form the holomorphic expression
\begin{align}
 E_N(x,t)
 &=N\int_x^{b_N}f_N(u,t)\,\dd u
       +\frac{f_N(x,t)-f_N(b_N,t)}2
 \notag\\
 &\quad+\sum_{k=1}^{K_N}
 \frac{B_{2k}}{(2k)!}N^{1-2k}
 \bigl[\partial_x^{2k-1}f_N(b_N,t)
       -\partial_x^{2k-1}f_N(x,t)\bigr].
 \label{eq:b1-optimal-EM}
\end{align}
The Bernoulli polynomial bound
$|B_{2k}(u)|\le4(2k)!/(2\pi)^{2k}$ together with
\eqref{eq:b1-Cauchy} bounds the Euler--Maclaurin remainder by
\[
 C N\left(\frac{K_N}{\pi aN}\right)^{2K_N}
 \le C e^{-cN}
\]
after decreasing $c_a$. This is uniform in the complex color disk and
on the physical high grid. Since the product of the exact ratios is
precisely $U_r$, it follows that
\begin{equation}\label{eq:b1-exponential-grid}
 U_{Nt,r}=e^{E_N(r/N,t)}[1+R_{N,r}(t)],
 \qquad |R_{N,r}(t)|\le C e^{-cN}.
\end{equation}
The empty terminal product equals one. The residual is retained as a
source term, rather than reinterpreted as a homogeneous solution.

The same explicit expression is holomorphic in the index $x$.
Using \eqref{eq:b1-ratio-first} in
\eqref{eq:b1-optimal-EM}, the higher Bernoulli terms sum to $O(N^{-1})$
on an inner fixed tube, and hence
\begin{equation}\label{eq:b1-interpolant-leading}
 e^{E_N(x,t)}=C(x,t)e^{NS(x,t)}[1+O(N^{-1})]
\end{equation}
uniformly holomorphically. Equations
\eqref{eq:b1-exponential-grid} and
\eqref{eq:b1-interpolant-leading} have distinct roles: the former is
an exponentially accurate comparison with the original finite source,
and the latter is a leading expansion of a specified complex function.

\subsection{Keeping the quadratic phase and determining the signed action}

The half-twist coefficient contains a quadratic exponential in the color. Such a phase is sometimes called a chirp. Its continuation is part of the finite formula, not an arbitrary choice of interpolation. We now combine it with the recurrence solution and the continued factorials; this calculation determines the local action and the nonzero amplitude used in the saddle.

Remove only the known external phase by defining
\[
 \ell_N^\eps(r,t)=e^{-\ii\eps\pi d/4}d_N^\eps(r,t),
 \qquad \eta_N=\exp\left(\frac{\ii\pi d}4
                               \sum_{\ell}\eps_\ell\right).
\]
The positive sign is the coefficientwise conjugate of the negative sign
on real parameters. Its holomorphic continuation uses Schwarz
reflection, not conjugation of a free holomorphic variable.

The original factorial phase gives the exact identity
\begin{equation}\label{eq:b1-color-phase}
 q^{h^2}\frac{D_{h,r}}{D_{0,r}}
 =\frac{G_N(r)^2}{G_N(r-h)G_N(r+h)}.
\end{equation}
It follows immediately from \eqref{eq:b1-marked-factorial}; the linear
and quadratic exponents in the two shifted arguments must both be kept.
Combining \eqref{eq:b1-signed-finite}--\eqref{eq:b1-D-source} with
\eqref{eq:b1-color-phase} gives the negative-sign interpolant
\begin{align}
 A_N^{\rm hol}(x,t)
 &=\frac{N^2G_N(2Nx)}
 {G_N(Nx)^2G_N(N(x-t))G_N(N(x+t))},
 \notag\\
 \nu_N(z)&=e^{-\pi\ii z+\pi\ii z(z+1)/N},
 \qquad P_N^{\rm ch}(x)=\frac{\kappa^{-1}e^{\pi\ii d/4}}{\nu_N(Nx)},
 \notag\\
 I_N^-(x,t)&=A_N^{\rm hol}(x,t)P_N^{\rm ch}(x)e^{E_N(x,t)}.
 \label{eq:b1-signed-interpolant}
\end{align}
The superscript on $P_N^{\rm ch}$ distinguishes this chirp factor from
the Casimir polynomial. At every integer $r$,
$\nu_N(r)=(-1)^rq^{r(r+1)}$. Its displayed analytic continuation is
fixed before any contour deformation. In particular,
\begin{equation}\label{eq:b1-chirp-cancellation}
 P_N^{\rm ch}(x)e^{NS(x,0)}
   =e^{\pi\ii(3/4-x-1/(2N))}.
\end{equation}
Changing the continuation of $(-1)^r$ would change a lattice frequency;
that change is not permitted. Equation
\eqref{eq:b1-exponential-grid} proves that
$I_N^-(r/N,t)$ agrees with the actual $\ell_N^-(r,t)$ up to a
relative exponential residual. For the positive sign set
$I_N^+(x,t)=\overline{I_N^-(\bar x,\bar t)}$; for the pure sign use
$I_N^0=A_N^{\rm hol}$, which agrees exactly on the grid.

Define
\begin{align}
 g(x,t)&=J(2x)-2J(x)-J(x-t)-J(x+t),
 \notag\\
 f_0(x,t)&=g(x,t),\qquad
 f_-(x,t)=g(x,t)+S(x,t)-S(x,0),
 \notag\\
 f_+(x,t)&=\overline{f_-(\bar x,\bar t)}.
 \label{eq:b1-local-actions}
\end{align}
The factorial expansion and \eqref{eq:b1-chirp-cancellation} yield
\begin{equation}\label{eq:b1-local-asymptotic}
 I_N^\eps(x,t)=N^{1/2}b_\eps(x,t)e^{Nf_\eps(x,t)}
                                      [1+O(N^{-1})],
\end{equation}
uniformly holomorphically on fixed smaller domains, with
\begin{align}
 b_0(x,t)&=
 \frac{\sqrt{\sin(2\pi x)}}
 {2^{3/2}\sin(\pi x)
       \sqrt{\sin(\pi(x-t))\sin(\pi(x+t))}},
 \notag\\
 b_-(x,t)&=b_0(x,t)e^{\pi\ii(3/4-x)}C(x,t),\qquad
 b_+(x,t)=\overline{b_-(\bar x,\bar t)}.
 \label{eq:b1-local-amplitudes}
\end{align}
All square roots continue their positive complete factorial branches.
At completeness,
\begin{equation}\label{eq:b1-amplitude-at-zero}
 \frac{b_\eps(x,0)}{b_0(x,0)}
 =\sqrt{\sin(\pi x)}e^{-\ii\eps\pi x/2},
 \qquad \eps\in\{-1,1\}.
\end{equation}
These factors are nonzero. The remaining $e^{-\pi\ii/(2N)}$ in
\eqref{eq:b1-chirp-cancellation} belongs to the stated inverse-power
error; it is not an omitted constant or parity phase.

\subsection{An exact source action identity and its Hessian}

The action above can also be characterized without integrating the
quadratic root. For each physical sign $\eps\in\{0,1,-1\}$ solve
near $s=t=0$ the equation
\begin{equation}\label{eq:b1-shear}
 t=s-\frac{\eps}{2\pi\ii}g_s(x,s).
\end{equation}
Its $s$ derivative at zero is $1-\eps\ii\cot(\pi x)$, which is
nonzero. Then the source-selected action satisfies
\begin{equation}\label{eq:b1-shear-action}
 f_\eps(x,t)=g(x,s)-\frac{\eps}{4\pi\ii}g_s(x,s)^2,
 \qquad (f_\eps)_x=g_x(x,s),\qquad (f_\eps)_t=g_s(x,s).
\end{equation}
The derivative identities follow by differentiating
\eqref{eq:b1-shear}. To identify the value with the actual selected
source, consider the negative sign. Put
$E=e^{2\pi\ii x}$, $H=e^{2\pi\ii t}$ and $T=e^{2\pi\ii s}$.
Exponentiating \eqref{eq:b1-shear} gives
\[
 H=\frac{T(E-T)}{ET-1},\qquad
 T^2+E(H-1)T-H=0,
 \qquad T=1\text{ at }H=1.
\]
For $K=e^{-g_x(x,s)}=-(1-E/T)(1-ET)/(1+E)^2$, elimination of $T$
gives
\begin{equation}\label{eq:b1-action-polynomial}
 K^2(1+E)^2+K[(H+H^{-1})E^2+2]+(1-E)^2=0.
\end{equation}
On the other hand \eqref{eq:b1-local-actions} and the selected quadratic
give $e^{-(f_-)_x}=E\mu D_t/(1+E)^2$, satisfying exactly
\eqref{eq:b1-action-polynomial}. At $H=1$, both constructions take
the branch $K=\tan^2(\pi x)$; their logarithms agree by their complete
marks. Hence their $x$ derivatives agree. The possible remaining function
of $t$ is zero by the limit $x\to1/2$: there $S=0$,
$J(1/2+t)=-J(1/2-t)$, and the implicit shear has $s=t$, $g_s=0$.
Schwarz reflection proves the positive sign. This fixes the source
action, not merely a Hessian with the same entries.

In particular, all local actions are even in $t$, equal $g(x,0)$ at
zero, and
\begin{equation}\label{eq:b1-local-hessian}
 (f_\eps)_{tt}(x,0)
 =-\frac{2\pi\cot(\pi x)}{1-\eps\ii\cot(\pi x)},
 \qquad \eps\in\{0,1,-1\}.
\end{equation}
Their real parts are strictly negative on the fixed high interval.

\subsection{The exact finite wire window and its lattice saddle}
\label{app:lattice-laplace}

We now specify the function to which the complex saddle argument is
applied. Let the fixed completed-root expression have $H$ leaves and
$K=2H-3$ primitive free integer wires. The last forced product is removed
using its exact relation $i+j=d$, whose coefficient is one. Every
remaining wire has the affine form
\[
 r_w=NL_w(y)+\delta_w,
 \qquad y=k/N,\quad k\in\ZZ^K,
\]
where the integer offsets $\delta_w$ are fixed by the original compiler.
Take a small fixed rectangular window $B$ about the strict complete
optimizer. All leaf coordinates lie in intervals of the form
\eqref{eq:b1-high-interval}, and every retained gate inequality is
strict on a fixed neighborhood of the closure of $B$.

The exact finite window is
\begin{equation}\label{eq:b1-exact-window}
 \mathcal C_N(\mathbf t)
 =\sum_{k/N\in B}
  \prod_{\ell=1}^{H}\ell_N^{\eps_\ell}(r_\ell,t_\ell)
  \prod_{g=1}^{H-2}\mu_N(r_{g,1},r_{g,2},r_{g,3}).
\end{equation}
Its support is independent of the complex colors. It is a finite,
separately even holomorphic function. For all sufficiently large $N$,
every state in this window is a true core state at every physical small
color deficit. The factors $\mu_N$ are the original positive Newton
coefficients, not their leading asymptotic replacements.

Replace each exact local factor by its canonical holomorphic
interpolant only after keeping its residual. The gate interpolation is
the same seven-factorial expression with $G_N$, and agrees exactly on
its integer grid. For the signed leaves, use
\eqref{eq:b1-signed-interpolant}. The total error in the finite sum is
bounded absolutely by
\begin{equation}\label{eq:b1-grid-error-before-saddle}
 \poly(N)\exp\{N(F_0+C\|\mathbf t\|^2-c)\}
\end{equation}
for some fixed $c>0$. Indeed there are polynomially many states, all
actions are even analytic perturbations of the common complete real
action, and the relative leaf errors in
\eqref{eq:b1-exponential-grid} are exponentially small. This estimate
is made before any oscillatory cancellation. Shrink the fixed color
polydisk so that \eqref{eq:b1-grid-error-before-saddle} is exponentially
below the nearby selected critical action.

The interpolated summand has phase
\begin{equation}\label{eq:b1-wire-phase}
 \mathcal F(y,\mathbf t)
 =\sum_{\ell}f_{\eps_\ell}(L_\ell(y),t_\ell)
  +\sum_g m(L_{g,1}(y),L_{g,2}(y),L_{g,3}(y)).
\end{equation}
At zero color it is the original complete action. Its wire Hessian at
the strict optimizer is negative definite. Thus it has one holomorphic
critical branch $y_*(\mathbf t)$, with
\[
 \Phi(\mathbf t)=\mathcal F(y_*(\mathbf t),\mathbf t).
\]
The local gate amplitude is
\[
 b_\mu(x,y,z)=2^{-3/2}
 \sqrt{\frac{\sin\pi(x+y)\sin\pi z}
 {\sin\pi x\sin\pi y\sin\pi(z-x)\sin\pi(z-y)\sin\pi(x+y-z)}}.
\]
Multiply all leaf and gate amplitudes by the exact leading offset factor
\begin{equation}\label{eq:b1-offset-factor}
 \exp\left[
  \sum_\ell\delta_\ell\partial_x f_{\eps_\ell}
  +\sum_g\nabla m(L_g(y))\mathbin{\cdot}\delta_g\right]
\end{equation}
to obtain the nonzero holomorphic density $b_E(y,\mathbf t)$.
The changed $f_\eps$ is used in \eqref{eq:b1-offset-factor}, including
at nonzero complex color; a pure offset factor would be incorrect.
The interpolated summand is
\[
 N^{3-H}b_E(y,\mathbf t)e^{N\mathcal F(y,\mathbf t)}
 [1+O(N^{-1})].
\]

For clarity we give the finite-lattice contour step. Choose the vertical
faces of the rectangle at half-grid values. In each coordinate the exact
summation kernel is $\pi N\cot(\pi Nz)$, with residue one at $z=k/N$.
On a vertical half-grid line, $|\cot(\pi Nz)|\le1$. Those faces lie
away from the strict real optimizer and retain a fixed phase gap.
On a horizontal line of fixed height $a>0$, the cotangent differs from
its constant value $\mp\ii$ by $O(e^{-2\pi Na})$. At complete color,
the real part of the phase grows under this small imaginary displacement
by at most $Ca^2$; the even color perturbation costs at most
$C\|\mathbf t\|^2$. Choose $a$ so $Ca^2<\pi a$, and then choose the
color polydisk small enough. Every nonconstant kernel term has an
exponential gap, including all side connectors. Applying this argument
successively in the $K$ primitive coordinates retains the full product
contour and its orientations. The constant horizontal terms give
$N^K$ times the ordinary integral. Replacing the moving half-grid
rectangle by $B$ costs only exponentially small boundary strips.

Translate that integral to the nearby complex critical point. The
endpoint connectors remain below its action, and the real tangent
Gaussian is uniformly nondegenerate. Extend its local Gaussian range
with exponentially small error. The odd first correction integrates to
zero, giving a relative $O(N^{-1})$ expansion. Thus
\begin{equation}\label{eq:b1-window-asymptotic}
 \mathcal C_N(\mathbf t)
 =N^{3/2}a_E(\mathbf t)e^{N\Phi(\mathbf t)}[1+O(N^{-1})],
 \qquad
 a_E(\mathbf t)=\frac{(2\pi)^{K/2}b_E(y_*(\mathbf t),\mathbf t)}
                    {\sqrt{\det[-\mathcal F_{yy}(y_*(\mathbf t),\mathbf t)]}}.
\end{equation}
The determinant square root is selected from the positive complete
Gaussian. The power is
$3-H+K/2=3/2$. The exponential grid error
\eqref{eq:b1-grid-error-before-saddle} is smaller than this answer by
the prior choice of color disk, and is added back only now. Both level
parities have the same primitive density and displayed normalization.

At completeness,
\begin{equation}\label{eq:b1-global-amplitude}
 \Phi(0)=F_0,\qquad
 a_E(0)=a_E^{\rm pure}(0)
 \prod_{\eps_\ell\ne0}
   \sqrt{\sin(\pi x_{\ell,*})}e^{-\ii\eps_\ell\pi x_{\ell,*}/2}\ne0.
\end{equation}
Evenness gives $\mathcal F_{yt_\ell}(y,0)=0$, so the wire Schur
correction to the color Hessian vanishes. With
$c_\ell=\cot(\pi x_{\ell,*})$, one obtains
\begin{equation}\label{eq:b1-global-hessian}
 \partial_{t_\ell}\partial_{t_j}\Phi(0)
 =-\frac{2\pi c_\ell}{1-\eps_\ell\ii c_\ell}\delta_{\ell j}
 =2\pi\ii\tau_\ell\delta_{\ell j},
 \qquad \tau_\ell=\frac{\ii c_\ell}{1-\eps_\ell\ii c_\ell},
 \qquad \Im\tau_\ell>0.
\end{equation}
On a smaller fixed complex box,
\begin{equation}\label{eq:b1-finite-packet-envelope}
 |\mathcal C_N(u+\ii v)|
 \le\poly(N)e^{N(F_0-a\|u\|^2+B\|v\|^2)},\qquad a>0.
\end{equation}
This bound is for the exact finite function
\eqref{eq:b1-exact-window}, by
\eqref{eq:b1-window-asymptotic}, not just its leading formal action.

\subsection{Returning from the local sum to the original evaluation}

We finish by stating precisely how the independent alias estimates are
used. The compositional proof in the main text gives a uniform positive
gap for every nonempty alias sector on the entire physical color cube.
It also gives a concave all-core majorant with the unique strict complete
maximizer used above. The latter uniqueness is global: it agrees with the
strict complete action near that point, and a second maximum of a concave
function would give a constant segment through that strictly concave
neighborhood.

Hence all core states outside the fixed wire window have a fixed gap,
uniformly in physical colors. Within that window the actual local
far-color recurrence estimate loses a fixed positive amount whenever
one physical deficit exceeds a fixed small corridor. Polynomial state
counting and the finite number of sectors therefore give constants
$\delta,\gamma>0$ such that the literal mixed scalar $\mathcal A_N$ obeys
\begin{align}
 \max_{0\le h_\ell\le\delta N}
 \left|\mathcal A_N(\mathbf h)-
          \eta_N\mathcal C_N(\mathbf h/N)\right|
 &\le\poly(N)e^{N(F_0-\gamma)},
 \label{eq:b1-source-restoration}\\
 \max_{\exists\ell:\,h_\ell>\delta N}
 |\mathcal A_N(\mathbf h)|
 &\le\poly(N)e^{N(F_0-\gamma)}.
 \label{eq:b1-color-complement}
\end{align}
The maximum in the second line is over the entire original physical
color cube, not just the low-core band. No alias is folded into the even
analytic packet. The external $\eta_N$ is exactly the retained unit
phase defined above.

The global coherent-core input here is its uniform exponential-rate
envelope, of the form $\exp\{N(\widetilde g+o(1))\}$ with the
$o(1)$ uniform over all physical indices and colors. It is not a
polynomial-prefactor assertion through every spectral collision.
To obtain \eqref{eq:b1-source-restoration}--\eqref{eq:b1-color-complement},
first fix the strictly positive variational and far-color margins, then
absorb the uniform $o(N)$ errors into, for example, half their minimum.
Polynomial counting and the remaining finite prefactors are harmless
after that reduction of $\gamma$. The polynomial-prefactor local packet
proved above is a separate estimate on its fixed regular high interval.

These estimates preserve the whole-source remainder through the finite
filling sum. For each fixed signed filling tuple, the exact Hermite
coefficient bound is $\sum_j|\beta_{j-1}^{-p_\ell}|\le2N(1+|p_\ell|)$.
Thus \eqref{eq:b1-source-restoration} and
\eqref{eq:b1-color-complement} incur only a polynomial multiplier in
that whole sum. The finite analytic packet, rather than its pointwise
leading approximation, is then used for the Fourier endpoint and charge
arguments. Its exact evenness, complex estimate
\eqref{eq:b1-finite-packet-envelope}, nonzero amplitude and Hessian
\eqref{eq:b1-global-hessian} are the needed inputs. The retained filling
analysis supplies the smaller filled action and a nonzero leading
coefficient; the source gap is not itself a replacement for that analysis.

The argument in this section requires
\eqref{eq:b1-high-interval}. No continuation through a quarter-index
dominance boundary, no additional root-leaf closure, and no theorem for
arbitrary tangle syntax is used in the manuscript's completed-root result.
